\subsection{3D Lensing Power Spectrum}\label{subsec:wl_power_spectrum}

This section develops the 3D power spectrum of weak lensing observables in the SFB framework.

\subsubsection{SFB decomposition of convergence}

When sources have extended redshift distribution, decompose in SFB basis:

\begin{equation}
\kappa(\mathbf{r}) = \sqrt{\frac{2}{\pi}} \int_0^\infty dk \sum_{\ell m} \kappa_{\ell m}(k) k j_\ell(kr) Y_{\ell m}(\theta, \phi)
\end{equation}

where $r$ is comoving distance.

\subsubsection{3D power spectrum definition}

\begin{equation}
\langle \kappa_{\ell m}(k) \kappa^*_{\ell' m'}(k') \rangle = C_\ell^\kappa(k,k') \delta_{\ell\ell'} \delta_{mm'}
\end{equation}

$C_\ell^\kappa(k,k')$ is a 2D matrix in wavenumber space at each $\ell$.

Under isotropy, $C_\ell(k,k') = P(k) \delta(k-k')$ (single function).

\subsubsection{Window function approach}

With finite survey selection $\phi(r)$:

\begin{equation}
C^{\text{obs}}_\ell(k_1, k_2) = \left(\frac{2}{\pi}\right)^2 \int_0^\infty dk \, k^2 P(k) W_\ell(k_1, k) W_\ell(k_2, k)
\end{equation}

where the window function encodes both the selection and the lensing kernel.

\remark{
This window function formalism converts complex 3D integrals into 1D convolutions, making numerical evaluation tractable.
}

See \cref{ch:observables} for how these 3D spectra project to observable angular power spectra.
